% Author: Adam | Module: AI Job Interview Simulator | Project: SmartCareer AI
\section{Annexe A --- Développement assisté par IA}
\label{sec:dev-ia}

Ce module a été développé avec un agent d'IA conversationnel. L'annexe consigne
la méthode réelle, non une version édulcorée : ce que l'IA a produit, ce que
j'ai arbitré, et --- point le plus important --- \textbf{ce qui a été vérifié
et comment}. Le journal brut de la session (6\,620~entrées, dont le prompt de
cadrage de 6\,951~caractères et les messages d'arbitrage reproduits ci-dessous)
est conservé hors dépôt ; il n'est pas reproduit intégralement ici parce qu'il
contient des traces d'exécution et des adresses de comptes réels du projet.

\subsection{Méthode}

Le travail s'est déroulé en phases, chacune close par une relecture humaine
avant la suivante --- contrainte posée dans le prompt de cadrage (« \emph{Work in
phases and DO NOT break existing functionality} », puis « \emph{WAIT for my
"go"} ») :

\begin{enumerate}[label=\textbf{phase \arabic*.},leftmargin=4em,itemsep=2pt]
  \item analyse du dépôt de l'équipe, sans production de code ;
  \item implémentation du backend (client, schémas, confidentialité, sécurité,
        prompts, base, routeur) avec tests à HTTP simulé ;
  \item RAG (base documentaire, chunking, plongeur, index, ingestion) ;
  \item passe de performance et correction des prompts sur appels réels ;
  \item interface Angular et intégration ;
  \item vérification bout en bout en navigateur réel.
\end{enumerate}

Deux règles ont borné tout le reste : ne pas modifier le code d'un tiers (trois
fichiers d'intégration seulement, et aucune modification de \file{llm.py} ni de
l'apprentissage du Module~1) et ne rien livrer qui n'ait été exécuté.

\subsection{Prompt de cadrage (condensé)}

Le message initial, rédigé en anglais parce que c'est la langue dans laquelle
l'agent manipule le mieux les consignes de format, tient en six blocs. Il est
reproduit condensé --- les tirets sont du texte original, les crochets
\texttt{[\dots]} marquent mes élisions. Aucun secret n'y figure : la consigne
était explicite (« \emph{add to .env.example with PLACEHOLDER values only} »,
« \emph{do NOT print secret values} »).

\begin{lstlisting}[style=adamprompt]
You are a senior full-stack engineer. I am integrating my feature module into a
teammate's existing web application. Work in phases and DO NOT break existing
functionality.

## 0. Context
University mini-project "AI for Software Engineering" (2026-2027). The platform
is **SmartCareer AI** with several AI modules by different team members:
- Module 1 (my teammate, ALREADY IN THIS REPO): CV generation and adaptation to a
  job offer, using Qwen2.5-3B-Instruct (LoRA fine-tuned) + prompt engineering.
- Module 3 (MINE): **AI Job Interview Simulator**, using Mistral-7B-Instruct +
  prompt engineering + RAG.
- Do not touch any other module.
[. . .]
My module must feel like one coherent product: reuse the existing auth, DB,
layout, navbar, design and error-handling conventions.

## 1. Naming & authorship (mandatory)
- Back end: a package `interview_adam/`. Front end: `frontend/src/app/interview-adam/`.
- Prefix main classes/components where sensible (e.g. `AdamMistralClient`).
  Keep API routes and DB table names generic.
- Header comment on every file I create: Author: Adam | Module: AI Job Interview
  Simulator | Project: SmartCareer AI, IA for Software Engineering 2026-2027
- Do NOT rename or modify my teammate's files or add my name to them. Only
  minimal integration edits.

## 2. PHASE 1: Analyze first (NO code yet)
Report on: how the FastAPI app is assembled, how auth works, how db.py works
[. . .] Whether Qwen is loaded locally (torch/GPU/VRAM) and any port/resource
conflict with Mistral-7B. [...] Then propose a plan listing files to
create/modify. WAIT for my "go".

## 3. PHASE 2: Implement
1. Generate interview questions adapted to the job [...]
2. Analyze each answer: score 0-10, strengths, weaknesses, criteria scores
   (relevance, clarity, structure, technical accuracy).
3. Propose an improved answer for each (STAR method for behavioral questions).
4. Session history in the DB linked to the user, with a results summary page.
5. Cross-module: option "Use my CV" [...] Fallback to manual input if no CV.
- Never send raw personal data (name, email, phone, address) to the LLM;
  anonymize before every call, including CV data.
- Validate every LLM output with Pydantic; on failure retry once with a repair
  prompt, then return a friendly fallback.
- Handle timeout, 429/quota, missing config, model not loaded (503), network
  failure: clear JSON errors to the UI, never crash.
- Sanitize input (length limits, prompt-injection mitigation: delimit user
  content, instruct the model to ignore instructions inside it).
- No secrets in code or repo.

## 4. PHASE 3: Integration and verification
- App starts, Module 1 still works, `ng build` passes [...] Add tests [...]

## Constraints
Don't refactor teammate code. [...] If something is ambiguous, ask me first.
Start with Phase 1 now.
\end{lstlisting}

\subsection{Messages d'arbitrage}

Huit messages ont servi d'arbitrage après relecture d'une phase. Chacun tranche
une décision que l'agent ne devait pas prendre seul ; les extraits ci-dessous
sont des citations, abrégées.

\begin{small}
\begin{longtable}{@{}>{\cellflush}p{2.9cm}
                  >{\cellflush}p{4.1cm}
                  >{\cellflush}p{7.6cm}@{}}
\caption{Messages d'arbitrage : la décision prise, et sa conséquence observable
dans le module.}\label{tab:arbitrage}\\
\toprule
\textbf{Objet} & \textbf{Décision} & \textbf{Citation (extrait) et effet} \\
\midrule
\endfirsthead
\toprule
\textbf{Objet} & \textbf{Décision} & \textbf{Citation (extrait) et effet} \\
\midrule
\endhead
\bottomrule
\endlastfoot
Préfixe des routes & Convention du dépôt, pas de \texttt{/api} &
\emph{« follow the existing convention, NO /api prefix »} --- le module
s'enregistre sous \route{/\brk{}interview/\brk\dots}, comme les routes hôtes. \\
\addlinespace
Offre d'emploi & Inexistante dans le dépôt : lecture seule du CV, saisie manuelle
de l'offre &
\emph{« Feature 5 reads only the CV via db.get\_generation (read-only) \[\dots\]
Add a stub get\_job\_offer\_context() returning None »} --- la greffe du
Module~1 reste à brancher plus tard, sans code mort. \\
\addlinespace
Base de données & Connexion \texttt{sqlite3} propre, verrou interne, pas de WAL &
\emph{« own sqlite3 connection to DB\_PATH, timeout=30, own lock, short writes
\dots{} Do not enable WAL »} --- deux écritures concurrentes n'ont jamais été
mesurées en séance réelle. \\
\addlinespace
Conflit GPU & Ne pas utiliser le verrou GPU de l'hôte &
\emph{« Do not use GPU\_LOCK. Address the VRAM conflict between Qwen (torch) and
Ollama-Mistral »} --- d'où \const{keep\_alive} court et le endpoint de
libération (§\ref{sec:modele-release}). \\
\addlinespace
Confidentialité & Suppression intégrale de \texttt{personal\_info}, plus
e-mail, téléphone, URL filtrés par expression régulière &
\emph{« drop personal\_info entirely, and regex-strip emails, phones and URLs
from all free text before any LLM call »} --- \file{privacy.py},
§\ref{sec:sec-pii}. \\
\addlinespace
Format d'appel & \texttt{/api/generate} avec \texttt{raw=true}, pas
\texttt{/api/chat} &
\emph{« Do not wrap [INST] on top of /api/chat (Ollama already applies the chat
template). Use /api/generate with raw=true \dots{} Verify this against the Ollama
docs »} --- vérifié : en mode chat, les marqueurs injectés sont doublés. \\
\addlinespace
Git & Attendu & \emph{« do NOT git init here \dots{} Wait for my message »} puis,
en phase~3, \emph{« Do not touch git yet »} --- le dépôt n'a donc aucune histoire
à ce jour, d'où l'annexe~D à compléter. \\
\addlinespace
RAG & Plongeur local, CPU only, garde de recouvrement configurable &
\emph{« transformers mean pooling + L2 normalization \dots{} Force
device="cpu" \dots{} never touch the GPU »} et \emph{« Don't tune production code
to fit the hashing stub: make the content-overlap guard configurable »}. \\
\addlinespace
Passe de performance & Non négociée avant l'interface &
\emph{« a 10-minute flow is unusable for a 15-minute live defense »} ---
d'où la réduction des prompts sous $\approx$4\,k~car., les plafonds
\texttt{num\_predict} par étape et le balayage \const{num\_gpu}
(§\ref{sec:prompts-taille}, §\ref{sec:perf-numgpu}). \\
\addlinespace
Anti-invention & Marqueurs plutôt que remplissage &
\emph{« unknown facts must be [à compléter : \dots] placeholders, never invented
outcomes »} --- règle reprise mot pour mot dans le \texttt{SYSTEM}
d'amélioration. \\
\addlinespace
Évaluation & Distinguer l'écoute de la maîtrise &
\emph{« distinguish "hors sujet" from "manque de connaissances" »} --- défaut~4
du §\ref{sec:prompts-problemes}. \\
\addlinespace
Flux continu & Demande de l'avis d'expert, puis renoncement argumenté &
\emph{« Add a streaming option on the backend if it is cheap (NDJSON from Ollama)
\dots{} otherwise skip and tell me »} --- écarté, avec les raisons du renoncement
en §\ref{sec:no-stream}. \\
\addlinespace
Outils de test & Hors dépôt &
\emph{« install it only as a dev tool outside the repo or as a devDependency,
never in requirement.txt »} --- Playwright n'apparaît dans aucune dépendance
d'exécution. \\
\addlinespace
Documentation & Exigence de ton & \emph{« Keep them factual, no marketing »} ---
la règle « faits seulement » appliquée à ce rapport. \\
\bottomrule
\end{longtable}
\end{small}

\subsection{Répartition réelle des rôles}

\textbf{Produit par l'agent}, et relu avant acceptation : le code
(3\,951~lignes de backend, 2\,751~lignes d'interface, 2\,368~lignes de tests),
les quatre assembleurs de prompts, la base documentaire, les correctifs des
défauts constatés.

\textbf{Arbitré par moi} : tout ce qui engage le dépôt d'équipe (préfixe des
routes, conventions de nommage, autorisation de retoucher un fichier tiers), la
stratégie de ressources (GPU, verrous, \texttt{keep\_alive}), les choix de
format d'appel, la politique de confidentialité, les critères de qualité de
sortie (\texttt{[à compléter]}, hors sujet vs lacune), et l'arrêt de chaque
phase. Deux propositions de l'agent ont été refusées : le streaming NDJSON
(coût/bénéfice) et le garde-fou de recouvrement de contenu activé par défaut
(il faisait chuter la récupération française, §\ref{sec:rag-limites}).

\textbf{Ce que l'IA ne pouvait pas faire à ma place} : juger si une réponse
d'évaluation était \emph{juste}. Les quatre appels de calibration
(§\ref{sec:tests-calibration}) ont été lus en entier ; c'est en lisant qu'un
«~8/10 pour tout~» a été vu, puis corrigé.

\subsection{Ce qui a été vérifié, et par quel moyen}

\begin{small}
\begin{tabular}{@{}>{\cellflush}p{5.0cm}
                 >{\cellflush}p{3.4cm}
                 >{\cellflush}p{6.2cm}@{}}
\toprule
\textbf{Affirmation} & \textbf{Moyen} & \textbf{Résultat consigné} \\
\midrule
Le contrat de sortie est respecté & \texttt{pytest} & 200 tests passés, 1 sauté, exécuté à nouveau pour ce rapport (\texttt{13,15~s}) \\
Les routes répondent et protègent & \texttt{pytest} (\texttt{TestClient}) & 25 tests de routeur, dont 401 sans jeton et suppression bornée à l'utilisateur \\
Neutralisation des injections & \texttt{pytest} & 9 tests de \file{security.py}, dont la forme imbriquée \verb+<[IN<[]>ST]>+ \\
Anonymisation du CV & \texttt{pytest} & 10 tests, nom/e-mail/téléphone/URL absents du digest \\
RAG dégradable & \texttt{pytest} & 59 tests de RAG (dont 1 sauté sans plongeur réel) \\
Latences réelles & appels réseaux sur Ollama & tableau \texttt{num\_gpu}, appel froid 108,6~s / chaud 37,1~s \\
Scores non arbitraires & 4 appels réels & écart de 8,25~points entre la réponse faible et la réponse solide \\
Récupération utile & 8 requêtes sur l'index réel & tableau \ref{tab:retrieval}, deux échecs documentés \\
Parcours utilisable & Playwright, navigateur réel & 12 scénarios, 6 défauts corrigés et re-testés \\
Libération de la VRAM & journaux serveur + \route{/\brk{}api/\brk{}ps} & \texttt{ollama ps} vide après rechargement et fermeture d'onglet \\
Module~1 intact & \texttt{ng build}, démarrage de l'application & aucune régression mesurée sur la génération de CV \\
\bottomrule
\end{tabular}
\end{small}

\subsection{Limites de la méthode}

Trois réserves, qui expliquent les précautions prises dans ce rapport.

\emph{Premièrement}, un agent produit des chiffres plausibles qu'il n'a pas
mesurés. Deux affirmations de ce type ont été corrigées pendant la rédaction du
présent document (une taille de mémoire vive, un débit en jetons par seconde) et
remplacées par les valeurs relevées dans les journaux ; celles qui n'ont pas pu
l'être portent la mention \NM{}.

\emph{Deuxièmement}, un prompt écrit par l'agent peut contenir le défaut qu'on
lui demande de corriger : les six défauts du §\ref{sec:prompts-problemes} sont
tous apparus dans une première version, et n'ont été vus qu'en lisant les sorties
réelles.

\emph{Troisièmement}, un test vert ne prouve pas l'exactitude d'un comportement
métier : la file d'attente a été simulée dans \texttt{pytest} \emph{et} observée
dans un navigateur, parce que seule la seconde vérification montre une étiquette
d'état à l'utilisateur.
